We refined the \edit{questions and workflow provided in the next
section} over a period of roughly \edit{two years},
incorporating \edit{many rounds of feedback}.

First, leveraging our own experiences as researchers with diverse
backgrounds working in different domains and institutions, we drew on
our knowledge of dataset characteristics, unintentional misuse,
unwanted \edit{societal} biases, and other issues to produce an initial set
of questions \edit{designed to elicit information about these topics. We}
then ``tested'' the\edit{se} questions by creating example datasheets for two
\edit{widely used} datasets: Labeled Faces in the Wild~\cite{lfw} and Pang
and Lee's polarity dataset~\cite{polarity}. \edit{We chose these
datasets in large part because their creators provided exemplary
documentation, allowing us to easily find the answers to many of the
questions.} While creating the\edit{se example} datasheets,
we \edit{found} gaps in \edit{the} questions, as well as redundancies
and lack of clarity. We \edit{therefore} refined \edit{the} questions
and distributed them to product teams in two major US-based technology
companies, in some cases helping \edit{teams} to create datasheets for
their datasets and observing where \edit{the} questions did not
achieve their intended objectives. Contemporaneously, we circulated an
initial draft of this paper to colleagues through social media and on
arXiv (draft posted 23 March 2018). Via these channels we received
extensive comments from dozens of researchers, practitioners,
and \edit{policy makers}. We also worked with \edit{a team of lawyers}
to review the questions from a legal perspective.

We incorporated this feedback to yield the questions and workflow
\edit{provided} in the next section\edit{:} We \edit{added and removed}
questions,
\edit{refined the content of the} questions\edit{,} and \edit{reordered the questions}
to better match the key stages of the dataset lifecycle. Based on our
experiences with product teams, we reworded the questions to
discourage yes/no answers, added a section on ``Uses\edit{,''} and
deleted a section on ``Legal and Ethical Considerations.'' We found
that product teams were more likely to answer questions about legal
and ethical considerations if they were integrated into sections about
the relevant stages of the dataset
\edit{lifecycle} rather than grouped together. \edit{Finally, f}ollowing feedback
from \edit{the team of lawyers}, we removed questions \edit{that} explicitly
ask\edit{ed} about compliance with regulations, and introduced factual
questions intended to elicit relevant information about compliance
without requiring dataset creators to make legal judgments.
